% Compiles on its own: latexmk "20 - hintergrund" (in this directory, see README.md).
% The subfiles package loads the preamble of the main file named here.
\documentclass[../main-german]{subfiles}
%\documentclass[../main-basic-listings-german]{subfiles}
%\documentclass[../main-german-thu]{subfiles}
%\documentclass[../main-german-university-of-hamburg]{subfiles}
\begin{document}
%%% ===============================================================================
\chapter{Background (SM1)}\label{sec:background}
%%% ===============================================================================
\progress{15}

\emph{Purpose and scope}.
The background chapter has multiple roles.
\begin{itemize}

  \item \emph{Preliminaries.}
        It has to provide all (and exactly the) information that is necessary to understand the methodological and technical parts of your work in the specific area of study.
        Assume as starting point another student in your degree who did not study the specific subject you are studying but has the task to understand your work.
        Which concepts, terms, definitions, etc. does the student have to know?
        Which formulas, symbols, etc. are standard in this topic?
        Only introduce definitions if you actually need them in any of the subsequent chapters.

  \item \emph{Related Work.}
        It has to provide a comprehensive discussion of all prior work in the area on this subject.
        Your discussion has to summarize these prior works and has to explain in which way the research question you are solving (\zcref{sec:introduction:research-question}) has not been solved yet because prior work had more limiting assumptions, addressed a different angle, their results are not complete etc. Depending on the subject you are studying, the related work part can be larger and warrant an entire chapter on its own, or be fully concluded within \zcref{sec:introduction:state-of-art}.

        You can close the related work discussion by clarifying the positioning and formulation of your research question (SM2) in relation to all the prior work, making more explicit whether you address an existing research question under different premises or whether you work on a modified or completely new research question.
\end{itemize}
\end{document}
